\section{Threats to Validity}\label{sec:threats}

The threats below are the limits of the design in Section~\ref{sec:methods} and of the calls in
Appendix~\ref{app:lookup}. They use no result beyond the runs already reported.

\begin{itemize}
\item \textbf{One processor model and one core.} Every time comes from one core of one \HostLArch{}
processor. Another microarchitecture may rank close cells differently, above all the cells at
$m=\SizeTen$ against gin.
\item \textbf{One compiler per language.} C and C++ were compiled by clang only, Rust by rustc and Go by
its own toolchain. GCC and MSVC builds were not timed. MSVC appears only in the costs of compile-time
tables.
\item \textbf{The native instruction set.} Every arm was built for L's own processor. A generic build
would give other code and other times, which may matter most for the small tables against gin.
\item \textbf{Loopback for H6.} The end-to-end run used loopback, one server core and one request in
flight per connection. A network, more cores or pipelining would change the share of routing.
\item \textbf{The minimal server.} It parses, routes and answers a fixed body. A full framework adds
work per request, which dilutes the difference between routers.
\item \textbf{External routers end to end.} The archived end-to-end run compares v2 with v1 in this
paper's server. It records no request rate for gin, matchit or any other external router.
Section~\ref{sec:e2e-gap} leaves that block out. A lookup result stays a statement about lookups.
\item \textbf{Development on the same host.} The engineering of v2 ran on L, on other seeds. The run's
seeds were held out, but the design choices may still suit L's microarchitecture better than others.
\item \textbf{The Drogon transcription.} Drogon's router needs a running application, so its lookup was
transcribed into the harness from its source. The file's hash and line ranges are pinned. A
transcription can still differ from the router as Drogon runs it.
\item \textbf{The gin shim.} gin's tree is vendored unchanged, but the parts of gin's engine that a
lookup uses are reproduced in a small shim. They are the reset of the context, the method's tree and
the search in gin's default configuration. chi's arm uses a similar shim of its multiplexer. A shim
can differ from the engine it stands for.
\item \textbf{The repeated ring.} Every timed pass repeats the same \RingSize{} queries. Predictors and
caches learn the ring, and a large table's hot part is at most the ring's routes. Real traffic, with
more distinct paths, may favour other structures. The Zipf and miss-heavy rings explore two departures
only.
\item \textbf{Different measurement boundaries.} C and C++ arms are called in process.
The Rust arms are called through one C function per lookup.
The Go arms run their timed loops inside Go, entered through cgo.
Their nulls remove the loops, and other costs of the boundary remain.
A ratio of a C++ time to a Go time crosses that boundary.
Appendix~\ref{app:lookup} lists the calls.
\item \textbf{The cut for slow arms.} A competitor whose passes ran on a prefix of the ring has
optimistic times and counters. None was the fastest competitor of an H1 cell, but the cut competitors
enter the median of H2's branch misses, which Table~\ref{tab:htwo} gives with and without them.
\end{itemize}
